\subsection{Casewise construction of the forced finite sets}
\label{subsec:new-casewise-witness-details}

We now record the finite set used in each combinatorial row, including the
open-endpoint issues.  This makes the Strategy~3 ownership independent of the
former reach-propagation chapter.

\subsubsection{One gap and no supercritical V role}

Normalize the only gap to
\[
 J=[B_0,1-A_1]=[\ell,r]\subset e_{0,1}.
\]
On every other edge choose an open-overlap point with coordinate $x_i$.
Because $A_i+B_i\le1$, one has
\[
 x_1\le r,
 \qquad
 x_i\le x_{i-1}\quad(2\le i\le5),
 \qquad
 x_5\ge\ell.
 \tag{6.72}
\]
The endpoint inequalities remain valid when a gap or overlap degenerates to a
singleton; the strict open choices are obtained by an arbitrarily small
perturbation and the final enclosure number is continuous.

Put $p=1-r$ and $q=\ell$.  For $T_1$, the two selected lower demands are
$p$ and $x_1\ge q$.  For $2\le i\le5$, they are
$1-x_{i-1}\ge p$ and $x_i\ge q$.  At $T_0$, they are
$1-x_5\ge p$ and $q$.  Thus every local role dominates $(p,q)$.

Let $c_*=c_{\max}(p,q)$.  The own role cannot contain
$(1-c_*)V_i$ by the definition of $c_{\max}$.  If an adjacent role is Vd0,
it has no positive interval on $r_i$.  If it is T3-like, its unique adjacent
capacity is $C_+$ or $C_-$ and is at most $c_*$ by
Lemma~\ref{lem:new-common-pair-domination}.  The other roles are excluded by
diameter.  Hence
\[
 (1-c_*)V_i\in U_C
 \qquad(0\le i\le5).
 \tag{6.73}
\]
The six points form a regular hexagon of circumradius $1-c_*$.  The convex
center role therefore contains its inball
\[
 \mathcal D_{h(1-c_*)}.
\]
The gap segment is exactly $J(p,q)$.  This proves the hypotheses of
Theorem~\ref{thm:new-complementary-gap} without any abstract residual-hull
argument.

\subsubsection{The T3-like neighboring capacity}

We verify directly that the T3-like supported endpoint is the radical branch
of Proposition~\ref{prop:new-neighbor-ray-formula}.  Use the translated
T3-like normal form with boundary reaches $a,b$ and shape parameter $D$.
The support-side relation is
\[
 R_0^2-DR_0+D^2=1,
\]
and the far endpoint on the forward neighboring ray is
\[
 u=1+a-\frac{R_0}{D}.
\]
The boundary reaches satisfy
\[
 a+b=\frac1D.
\]
Solving the quadratic for $R_0/D$ gives
\[
 \frac{R_0}{D}
 =\frac{1+\sqrt{4(a+b)^2-3}}2.
\]
Therefore
\[
 \boxed{
 u=a+\frac12-\sqrt{(a+b)^2-\frac34},
 }
 \tag{6.74}
\]
which is exactly the final radical branch of $C_+(a,b)$.  Thus no additional
T3-specific radial function occurs in the one-gap proof.

A T3-like role is nonsupercritical.  Hence its boundary sum is at most one,
and (6.74) is bounded by
\[
1-\min\{a,b\}\le c_{\max}(p,q)
\]
whenever $(a,b)$ dominates the common pair $(p,q)$.  This is the explicit
reason why one or two T3-like roles do not shrink the regular radial witness
hexagon in Proposition~\ref{prop:new-nplus-zero-gap-closures}.

\subsubsection{Two gaps}

For CE2, write the two center intervals as
\[
 I_L=\left[\frac{k}{W},R+\alpha\right],
 \qquad
 I_R=\left[\frac{k}{R},W+\delta\right].
\]
If both contain gaps, then the incident endpoint roles give
\[
 A_1\ge R-\delta=q,
 \qquad
 B_5\ge W-\alpha=p.
\]
Choose handoffs $x_1,\ldots,x_4$ on the four intervening edges.  The five
nonsupercritical roles give
\[
 p<x_4<x_3<x_2<x_1<1-q.
 \tag{6.75}
\]
Every one of $T_1,\ldots,T_5$ therefore dominates $(p,q)$.  In particular,
the two transverse witnesses
\[
 D_2=(1-c_{\max}(p,q))V_2,
 \qquad
 D_4=(1-c_{\max}(p,q))V_4
\]
are excluded from all V roles.  The center exits at distances
$\delta$ and $\alpha$ on these two rays.  Theorem
\ref{thm:new-ce2-short-ray} states that one witness lies strictly beyond its
exit.  This argument does not inspect $A_0+B_0$, so it simultaneously handles
$N_+=0$ and $N_+=1$.

\subsubsection{One gap, one supercritical Vd0 role}

In the all-Vd0 branch, define the six actual radial endpoints
\[
 P_i=(1-C_i)V_i.
\]
Because $U_i$ is open, the endpoint does not belong to $U_i$.  The adjacent
roles have no positive support on $r_i$, and nonlocal roles are separated by
distance.  Hence $P_i\in U_C$.

Suppose a second open unit triangle $U'$ contains
\[
 O,M_0,X(\ell),X(r),P_0,\ldots,P_5.
\]
The fixed six V roles together with $U'$ cover every radial arm: $U'$ contains
$[O,P_i]$, and $U_i$ contains a relatively open interval from $V_i$ past every
point short of $P_i$.  They also cover the perimeter, since the unique missing
boundary component is the gap and $U'$ contains its two endpoints and hence
its full segment.  Thus all direct midpoint and signed-center conclusions
apply to $U'$.

If $d_i'$ is the exit of $U'$ on $r_i$, containment of $P_i$ says
\[
 d_i'\ge1-C_i,
 \qquad
 C_i\ge1-d_i'.
 \tag{6.76}
\]
This is the only connection between the finite set and the local reach
calculation.  In CE2, (6.76) supplies radial reaches $1-\alpha$ at $T_4$ and
$1-\delta$ at $T_2$; the direct threshold alternatives of the proof of
Proposition~\ref{prop:new-nplus-one-all-vd0} contradict the terminal bound at
$T_0$.  In CE1, (6.76) supplies the reaches used by Proposition
\ref{prop:new-ce1-direct-certificate}.  No statement about a residual polygon
is substituted for these actual endpoint facts.

For a singleton gap, replace $r=\ell$ by $r_n>\ell$ with $r_n\downarrow\ell$.
The diameter inequalities and the unique-supercritical pattern persist for
large $n$.  The witness sets converge in the Hausdorff metric, and
$\Lambda$ is continuous by the support formula.  Hence the positive-length
proof passes to the singleton limit.

\subsubsection{The T3-like rescuer endpoint}

In Proposition~\ref{prop:new-one-t3-terminal}, the finite point
$P_T=(1-u)V_1$ is forced into the center for a stronger reason than a generic
radial capacity.  The T3-like open trace begins at the vertex side at $c$ and
ends at $u$, so its O-side endpoint is not in the open role.  The adjacent
supercritical role misses $M_1$ and therefore, by convexity, misses the entire
O-side segment.  All remaining roles are Vd0.  Thus the segment from $O$ to
$P_T$ is owned by the center.

The two inequalities (6.44) then compare the far boundary tail with the one
free strict-supercritical envelope.  They do not propagate a numerical state
through the four middle roles.  Instead, adding the four actual edge-cover
inequalities gives one path budget:
\[
(A_2+B_2)+(A_3+B_3)+(A_4+B_4)+(A_5+B_5)
\ge4+h_T-B_1>4.
\]
This direct sum is the full global contradiction.

\subsubsection{The Vd radial points}

In the adjacent Vd placement, define
\[
 q_2=1-\delta.
\]
The two roles capable of reaching the center interval on $r_2$ have exact
far endpoints bounded by
\[
 u_{1\to2}<q_2,
 \qquad
 C_2<q_2.
\]
Therefore any point of $r_2$ between the larger of these two endpoints and
$q_2$ is missed by all seven roles.  This is a literal one-point obstruction,
not a perimeter reformulation.

In the nonadjacent placement, the point
\[
 D_\tau=\min\{\rho_R,\rho_L\}V_\tau
\]
is beyond the Vd own-radial trace by (6.70), and beyond the center trace by
(6.49).  Its adjacent roles are Vd0 and its nonlocal roles are separated by
diameter.  This is again a literal uncovered point.

In the Vd1 rescue placement, the O-side endpoint of the supported interval
plays exactly the same role as $P_T$ in the T3-like case.  Only the local
normal-form proof of the two inequalities changes; the center-hiding and
four-role boundary sum are identical.

\subsubsection{The corrected replacement placement}

When neither distinguished role is based at $V_0$ and the exceptional role is
Vd1, the shared edge is center-free.  The corrected replacement constructs
one nonsupercritical Vd0 role in each of the two correct vertex charts.  It
preserves all six boundary and radial skeleton traces needed for coverage but
does not claim to preserve the number of gaps.  After replacement, recompute
$N'_{\rm gap}$.

If $N'_{\rm gap}=0$, the six nonsupercritical Vd0 roles cover the perimeter,
contradicting the strict boundary-overlap sum.  If $N'_{\rm gap}=1$, use the
common radial disk and complementary gap.  If $N'_{\rm gap}=2$, use the two
transverse short-ray witnesses.  Hence the replacement terminal has no return
to a reach-composition argument.
